\chapter{量子行为}

\section{原子力学}

“量子力学”描述物质和光的行为的各方面细节，特别是发生在原子尺度上的事件。在微小的尺度下事物的行为一点也不像我们有着直接经验的任何事物。它们的行为既不像波动，又不像粒子，也不像云雾，或弹子球，或悬挂在弹簧上的重物，总之不像我们曾经见过的任何东西。

牛顿认为，光是由微粒构成的，但是，之后发现光的行为像波动。然而，后来（在20世纪初叶）人们发现，光的行为有时确实又像粒子。又譬如，在历史上，电子起先被认为像粒子，后来发现它在许多方面的性质像波。所以，实际上它表现得两者都不像。现在我们不再说它到底是什么，我们说：“它什么都不像。”

然而，运气总算还好：电子的行为很像光。原子客体（电子、质子、中子、光子等等）的量子行为都是相同的，它们都是“粒子波”，或者随便什么你愿意称呼的名称。所以，我们所学的关于电子（我们将用它作为例子）的性质也可应用到所有的“粒子”，包括光子上。

在20世纪的前四分之一，有关原子与其他小尺度粒子行为的知识逐渐积累起来，给出了微小物体是如何活动的一些线索，由此也引起了越来越多的混乱，到1926和1927年，薛定谔、海森伯与玻恩终于解决了这些问题。他们最后对微小尺度物质的行为作出了协调一致的描述。本章中我们将开始研究这种描述的主要特点。

因为原子的行为与我们的日常经验不同，所以很难适应它，而且对每个人——不管是新手，还是有经验的物理学家——都显得奇特而神秘。甚至专家们也不能以他们所想要的方式去理解原子的行为，而且这是完全有道理的，因为一切人类的直接经验和所有的人类的直觉都只适用于大的物体。我们知道大物体的行为将是如何，但是在小尺度下事物的行为却并非如此。所以我们必须用一种抽象的或想象的方式，而不是把它与我们的直接经验联系起来的方式来学习它。

在本章中，我们将直接讨论以最陌生的方式出现的神秘行为的基本特征。我们选择用来考察的现象不可能，绝对不可能，以任何经典方式来解释，但它却包含了量子力学的核心。事实上，它包含着独一无二的奥秘。我们不能通过“说明”它如何作用来消除这个奥秘。我们只是告诉你，它是怎样起作用的。在告诉你它怎样起作用的同时，我们也将告诉你所有量子力学的基本特色。